\section{Remaining Work}
\label{sec:remaining}

The next step is a front-end demonstration of the measured capacity.
Several users each stream video to one server that runs \optd{}, and each user's browser renders that user's Gaussians with WebGPU, for example with WebSplatter~\cite{websplatter2026}.
A dashboard beside the rendered views shows the users served, each user's update latency and the GPU memory in use, from the readings the harness already records.
Adding users one at a time should expose the limit that the capacity sweep found on the test clip.
The demonstration also adds work the harness omits, serializing each update's Gaussians and sending them to a browser, and whether this lowers the capacity of \optd{} is the first question it answers.

Three extensions follow, in this order.
The first is the server's behavior beyond capacity.
Our harness admits every user and drops no update, so past capacity each user's queue grows for as long as the run lasts. % G2
Two and three users past the capacity of \optd{}, the worst user's p95 latency reached 3.6--8.5~s within a 60~s run. % G16
We plan to add two policies and measure each user's tail latency and the fairness of latency across users under overload.
The first is admission control that refuses a new user when the admitted users' declared update rates would exceed the measured throughput, as Clockwork rejects requests that would miss their deadline~\cite{clockwork2020}.
The second keeps only each user's latest keyframe, replacing a queued update that has not started, since it loses its value once a newer keyframe exists.
Fairness needs its own measurement because sharing modes with equal capacity can still give users different latency, as time-slicing and MPS did at five AnySplat users.

The second extension runs each user in its own container under Kubernetes.
The NVIDIA device plugin offers a time-slicing mode, documented as not isolating the workloads that share a GPU, and an experimental MPS mode that gives each replica an equal share of memory and compute~\cite{k8s_device_plugin}.
We found no published measurement of how well either mode isolates users, and plan to rerun the capacity sweep with one pod per user under each mode.

The third extension repeats the measurements, with the same configurations and capacity rule, on the H200 GPUs of the nano4 cluster.
The H200 also supports MIG, which gives each instance its own L2 cache and DRAM bandwidth and isolates faults between instances~\cite{nvidia_mig}, while MPS documents only a limited form of error containment between its clients~\cite{nvidia_mps}.
The RTX 6000 Ada does not support MIG~\cite{nvidia_mig}, and an earlier probe found MIG disabled on the nano4 node and in its scheduler, so a MIG configuration depends on the cluster's administrators enabling it. % not in LOG.md: nano4 probe job 682344, 2026-11-02, project memory note

\section{Conclusion}

We asked how many users one RTX 6000 Ada serves when each user sends one keyframe per second to a feed-forward 3DGS reconstruction server and every user must keep p95 latency within 1~s. % G6, G15, G16
Time-slicing and MPS served the same number of users in every setting we measured, and batching users' updates into one call lowered throughput. % G2, G3, G6, G7, G9
Reducing the work per update raised the capacity from 4 users with \serial{} to 10 with \opt{} and 14 with \optd{}, for 0.19~dB of keyframe PSNR on the DAVIS 2017 validation set against full precision. % G15, G16
At 14 users the card held 4,884~MiB, under 10\% of its memory, so memory was not the limit. % G16
The capacity results hold for one GPU model and one test clip that every user replays, with two repeats per point, and the batching comparison ran once per batch size. % G2, G3, G6, G7, G9, G15, G16

\section*{Availability}

The measurement harness, the experiment scripts and the figure script are in the project repository, and every script after the first two experiments carries its predictions, written before the run. % G1 and G2 have no runner header; G3-G16 runners do
The project's experiment log records each number in this report with the command that re-derives it from the run's output.
